% Part 2 results and tracking. Numbers come from a3/scripts/part2_wandb.py (see a3/part2_summary.csv).

\subsection*{Results}

We tracked every run in the team's Weights \& Biases (W\&B) project:
\begin{center}
\small\url{https://forge.coreweave.com/wandb/vraj-shah2121-university-of-toronto/nanochat/workspace}
\end{center}
The three full runs are \code{d12\_baseline}, \code{d12\_geglu} and \code{d12\_rope50}. Each logged validation bits per byte (BPB) every 250 steps, training loss, throughput and MFU every 100 steps, and the CORE score at steps 2000 and 2520. CORE was capped at 500 examples per task. Table~\ref{tab:part2-results} summarizes the runs, and Figure~\ref{fig:part2-training} compares their training curves.

The W\&B configs differ only in the flag under test (\code{--mlp-type=geglu} or \code{--rope-frac=0.5}) and the run name. All three runs use seed 42 and 2,520 steps, with the same estimated total of $1.0 \times 10^{18}$ FLOPs. All three start from a validation BPB of 3.16867, because the zero-initialized output projections make the changes invisible at step 0.

\begin{table}[H]
\centering
\caption{Picochat (d12) ablation results. For validation BPB, lower is better; for CORE, higher is better. Gaps are relative to the baseline. Throughput and MFU are medians of the logged steps after the first five.}
\label{tab:part2-results}
\begin{tabular}{@{}lrrr@{}}
\toprule
 & Baseline & GeGLU & Partial RoPE (50\%) \\
\midrule
Hardware & $8 \times$ H200 & $8 \times$ H200 & $8 \times$ H100 \\
Final validation BPB & 0.84537 & 0.84754 & 0.84505 \\
\quad Gap to baseline & -- & $+0.00217$ & $-0.00031$ \\
\quad Gap (\%) & -- & $+0.26\%$ & $-0.04\%$ \\
CORE at step 2000 & 0.1570 & 0.1547 & 0.1589 \\
Final CORE (step 2520) & 0.1734 & 0.1656 & 0.1705 \\
\quad Gap to baseline & -- & $-0.0078$ & $-0.0029$ \\
\quad Tasks better / worse (of 22) & -- & 9 / 13 & 8 / 13 \\
Training time (min) & 5.23 & 5.45 & 5.41 \\
Throughput (million tokens/s) & 4.33 & 4.16 & 4.10 \\
MFU (\%) & 41.6 & 39.9 & 39.4 \\
\bottomrule
\end{tabular}
\end{table}

\begin{figure}[H]
\centering
\includegraphics[width=\textwidth]{part2_training.pdf}
\caption{Training curves from the W\&B runs. The validation and training curves almost overlap, so the right panel plots each change's validation gap to the baseline. The plots start after the steep early drop (step 250 for validation, step 200 for training loss).}
\label{fig:part2-training}
\end{figure}

\subsection*{Impact of Each Change}

\textbf{Overall.} Neither change clearly improved the model at this scale, which fits the small or null effect we predicted in Part 1. Figure~\ref{fig:part2-core} shows why the CORE gaps are weak evidence: individual tasks move by as much as 0.10, in both directions. The standard error of the mean gap across tasks is about 0.008, so GeGLU's gap is 0.95 standard errors from zero and partial RoPE's is 0.34.

\textbf{GeGLU.} GeGLU has a slightly worse validation BPB than the baseline at all 11 evaluations after step 0, with a gap between 0.0016 and 0.0036 that had not closed by the end of training. Its CORE score is also lower, though within the task-level noise above. These evaluations come from one training run and are correlated, so the consistent sign is not 11 independent results. On the same GPU model as the baseline, GeGLU ran about 4\% slower in both throughput and training time. Its parameter count and estimated FLOPs match the baseline, so this slowdown is not captured by the FLOP estimate. We do not see the gain reported for gated MLPs in the literature \cite{glu}, but this run cannot rule one out: it is short and small, and it has only one seed.

\textbf{Partial RoPE (50\%).} We cannot distinguish partial RoPE from the baseline. It trails early (a BPB gap of $+0.0052$ at step 250), catches up by about step 1500, and finishes $0.00031$ ahead. It also led on CORE at step 2000 ($+0.0019$) but trailed at step 2520 ($-0.0029$), and that change of sign suggests the gap is mostly noise. ModernBERT reported a slight gain at 50--75\% \cite{modernbert}, and we neither confirm nor contradict that here. We also cannot read its throughput as a cost of the change: this run used H100s while the baseline used H200s, and the change rotates fewer dimensions, so it should not add work.

\textbf{Limitations.} We trained one model per configuration, so we have no direct measure of seed noise. The closest evidence is two identical 100-step smoke tests with the same seed, whose validation BPB differed by 0.00014. That shows only that runs are not fully deterministic, and it is a lower bound on noise, since changing the seed should vary results more. The GeGLU gap is about 15 times that figure and the partial RoPE gap about 2 times. The hardware also differs across runs, as noted above. It should not affect the loss curves much, but it does limit the throughput comparison to baseline versus GeGLU.

\begin{figure}[H]
\centering
\includegraphics[width=0.9\textwidth]{part2_core_tasks.pdf}
\caption{Final CORE gap to the baseline on each of the 22 tasks. The solid line marks zero and the dashed line marks the mean gap. Tasks are sorted by their average gap across the two changes. Both changes gain on some tasks and lose on others.}
\label{fig:part2-core}
\end{figure}
